\section{Open Questions}

The following five questions are, in the honest judgment of this replication, the most
important \emph{still-open} questions that the paper leaves unresolved and that a
follow-on study could productively attack. They are grounded in the re-read of Pariset
et al.\ 2020 and the specific artifacts that are missing (6/22 rule).

\begin{enumerate}

\item \textbf{Which specific genetic loci or DDR pathway variants drive the $\sim$2$\times$
spread in HZE relaxation time $\tau$ across the 15 inbred strains, and does the
$(\tau_{\text{HZE}}, q_{\text{HZE}})$ phenotype map to known Trp53/Atm/Brca1/53bp1/H2ax
variant classes?}

\emph{Basis.} Fig.~6 identifies MegaMUGA SNP peaks but the genotype file is not
deposited (data-blocked at replication level), and Table 1A per-particle Pearson matrix
is unreproducible. The paper leaves the mechanistic gap: strain-level $\tau$ is
observed, but the causal genetic driver is only speculated from SNP peaks near
candidate DDR genes.

\emph{Next steps.} Request MegaMUGA genotypes + per-cell foci counts from the Costes
lab under a data-use agreement; intersect strain-level $\tau_{\text{HZE}}$ rank order
with published Trp53/Atm/H2ax variant panels for these 15 strains (Mouse Phenome
Database); run eQTL-style regression of $\tau$ on cis-SNPs within a 500\,kb window of
the 8 core DDR genes; validate hits by CRISPR knock-in of the candidate variant into a
reference strain and re-measuring 53BP1 kinetics.

\item \textbf{Does the ex-vivo skin fibroblast 53BP1 repair phenotype actually predict
in-vivo tissue-specific tumorigenesis, or is the Fig.~7B / 7C signal an artifact of
small-$n$ over-fitting across 10 CC strains and 4 MTB strains?}

\emph{Basis.} Fig.~7B $r=0.61$, $n=10$ yields derived $p_{\text{two-sided}} \approx
0.061$ (borderline non-significant), and Fig.~7C at $n=4$ has an analytical critical
$|r|=0.950$ with zero organs surviving Bonferroni across 19 organs. The paper's
central translational claim (``ex-vivo predicts in-vivo susceptibility'') is therefore
statistically thin.

\emph{Next steps.} Enlarge the CC cohort to $\sim$30 strains (feasible via the UNC
Collaborative Cross resource) to bring Fig.~7B to $n \geq 20$; reduce the organ
multiple-comparison burden by pre-registering 3--4 organs of highest biological
interest; run leave-one-strain-out cross-validation; attempt an orthogonal readout
($\gamma$H2AX foci or comet assay) on the same strains to test whether \emph{any}
DSB kinetic marker predicts in-vivo cancer, not just 53BP1.

\item \textbf{How stable is the $(\tau, q, \text{RIF}_{\max})$ fit under alternative
time-sampling and noise-model assumptions --- do 4 time points at 0.5/2/4/24 h
post-irradiation actually identify all three parameters, or does the fit collapse to
an effective 2-parameter surface with $\text{RIF}_{\max}$ and $q$ trading off?}

\emph{Basis.} The original-pass identifiability MC (Claim 11) passed under the paper's
stated noise model, but the paper never reports the parameter covariance matrix or
profile-likelihoods. With only 4 non-uniform time points and biexponential decay $+$
residual plateau, $\text{RIF}_{\max}$ and $q$ are known to be near-degenerate for
$\tau$ in the 2--8\,h range (much of the observed strain spread). If real, the reported
strain-level $(\tau, q)$ may be aliased.

\emph{Next steps.} Recompute the Fisher information matrix at each strain's best-fit
$(\tau, q, \text{RIF}_{\max})$ and report the condition number; bootstrap over synthetic
4-point samples with the paper's stated Poisson foci-count noise; simulate a 6-time-point
protocol (adding 8\,h and 48\,h) and quantify how much parameter uncertainty shrinks;
publish this as a design-of-experiment recommendation for follow-up ex-vivo studies.

\item \textbf{Does the strain-level ordering of DDR kinetics translate to human
inter-individual variability, and can a 53BP1 foci assay on human primary fibroblasts
(LCLs or dermal punch biopsies) recover a comparable dynamic range predictive of
radiotherapy toxicity?}

\emph{Basis.} The paper's motivation is human-radiotherapy prediction, but it never
bridges from mouse strains to human individuals. The $\sim$2$\times$ $\tau$ range
across inbred mice is a \emph{floor} estimate for outbred humans, and it is unknown
whether the assay format transfers cleanly (dose calibration, foci-detection
thresholds, replication-stress background foci all differ between species).

\emph{Next steps.} Adopt the ex-vivo protocol on a panel of 30--50 human LCLs from the
1000 Genomes reference set spanning ancestry groups; benchmark against archived clinical
radiotherapy toxicity cohorts (REQUITE, RAPPER); cross-check with existing radiogenomic
GWAS hits for late toxicity (SNPs near TGFB1, XRCC1, ATM); publish a species-bridging
equivalence table before proposing the assay as a clinical stratifier.

\item \textbf{Given the near-complete lack of data deposition (no supplement, no
figshare/Zenodo/GitHub, no raw foci counts, no genotypes), what curation or
re-analysis of the underlying NASA Space Radiation Laboratory (NSRL) beamline
campaigns would be needed to make Pariset 2020 fully reproducible in the FAIR sense?}

\emph{Basis.} Re-pass full-text parse confirmed \emph{no} deposited data. This is
downstream of a multi-year NSRL campaign involving 76 mice $\times$ 15 strains
$\times$ 3 LETs $\times$ multiple doses --- an expensive, non-reproducible-from-scratch
dataset. Yet no per-cell foci table, no per-mouse metadata, no beamline dosimetry log
is published. This is the single largest reproducibility gap and blocks Table 1A
entirely.

\emph{Next steps.} Contact NSRL / NASA HRP data-archive coordinators to determine
whether the raw foci counts and per-mouse dosimetry \emph{are} archived internally
but not linked; if so, propose a re-release to a FAIR repository (Zenodo with
restricted-access DUA); if not, propose a limited re-collection focused on the 4
Table 2 strains that failed our quadrant match, to disambiguate paper-vs-digitization
noise; draft a companion ``data supplement'' paper with per-strain, per-particle
$(\tau, q, \text{RIF}_{\max})$ tables that the community can rebuild Table 1A from.

\end{enumerate}
